\chapter{Постановка задачи}
\label{ch:problem}

\section{Пространство поиска архитектур}
\label{sec:searchspace}

Каждый эксперт представляет собой свёрточную сеть, собранную из
\emph{ячейки} (cell) --- ориентированного ациклического графа, вершины
которого являются вычислительными узлами. Мы используем принятое в
реализации кодирование ячейки: ячейка содержит фиксированное число
узлов, и каждый узел применяет одну операцию из конечного множества
$\mathcal{O}$ (свёртки с разными размерами ядра, пулинг, тождественное
преобразование) к выходу одного из предшествующих узлов; первый узел
читает вход ячейки. Таким образом, архитектура полностью задаётся двумя
дискретными выборами для каждого узла --- \emph{какая} операция
применяется и \emph{в какой} узел она направляется, то есть какой из
предыдущих узлов подаётся ей на вход. Все эти выборы собираются в единый
вектор архитектуры $\boldsymbol{\alpha}$, так что $\boldsymbol{\alpha}$
описывает одну конкретную ячейку. Поскольку и множество $\mathcal{O}$, и
число узлов конечны, пространство поиска $\mathcal{A}$ всех допустимых
$\boldsymbol{\alpha}$ конечно, а поиск архитектуры сводится к выбору
наилучшего $\boldsymbol{\alpha} \in \mathcal{A}$. Для архитектуры
$\boldsymbol{\alpha}_k$ с обученными параметрами $\boldsymbol{\theta}_k$
соответствующую нейронную сеть обозначим
$f(\boldsymbol{\alpha}_k, \boldsymbol{\theta}_k): \mathcal{X} \to \mathcal{Y}$.

\section{Модель смеси экспертов}
\label{sec:moe}

Модель смеси экспертов объединяет $K$ таких сетей, каждая из которых
призвана специализироваться на своей области входного пространства.
Обозначив через
$\boldsymbol{\alpha} = [\boldsymbol{\alpha}_1^{\!\top}, \ldots, \boldsymbol{\alpha}_K^{\!\top}]^{\!\top}$
конкатенацию архитектур экспертов, прогноз MoE записывается как
взвешенная по маршрутизации комбинация выходов экспертов
\begin{equation}
\label{eq:moe_pred}
g(\boldsymbol{\alpha}, \boldsymbol{\theta}, \boldsymbol{w}_r)(\boldsymbol{x})
= \sum_{k=1}^{K} r_k(\boldsymbol{x}, \boldsymbol{w}_r)\, f(\boldsymbol{\alpha}_k, \boldsymbol{\theta}_k)(\boldsymbol{x}),
\end{equation}
где значения шлюза $r_k(\boldsymbol{x}, \boldsymbol{w}_r) \geq 0$,
$\sum_{k=1}^{K} r_k = 1$, вырабатываются маршрутизирующей сетью с
параметрами $\boldsymbol{w}_r$.

\section{Кластеризация данных и маршрутизация}
\label{sec:clustering}

Перед поиском архитектуры датасет разбивается на $M$ кластеров
$\mathcal{D} = \mathcal{C}_1 \cup \cdots \cup \mathcal{C}_M$, каждый из
которых группирует объекты схожей структуры (например, семантический
класс или модальность данных). Обозначим через
$\mathcal{C}_m \subseteq \{1, \ldots, N\}$ индексы объектов кластера
$m$, а через $m(n)$ --- номер кластера объекта $n$. Маршрутизация
задаётся на уровне кластеров: пусть $r_{mk} \in [0, 1]$ --- вероятность
приписывания кластера $m$ эксперту $k$, $\sum_{k=1}^{K} r_{mk} = 1$;
эти величины собираются в матрицу маршрутизации
$\mathbf{R} \in [0, 1]^{M \times K}$, а её столбец
$\mathbf{r}_k = \mathbf{R}_{:,k}$ кодирует кластеры, за которые отвечает
эксперт $k$. Итоговое жёсткое разбиение восстанавливается как
$I_k = \{m : \arg\max_{k'} r_{mk'} = k\}$.

В отличие от стандартного обучения MoE, где при фиксированной
архитектуре обучаются лишь параметры шлюза, здесь \emph{нахождение
маршрутизации $\mathbf{R}$ само является частью поиска архитектуры}:
наилучшая архитектура эксперта зависит от того, на каком подмножестве
кластеров он обучается, что связывает выбор архитектуры со
специализацией под данные.

\section{Целевая функция}
\label{sec:objective}

Поиск архитектуры MoE формулируется как задача максимизации
правдоподобия. Требуется выбрать для каждого эксперта $k$ архитектуру
$\boldsymbol{\alpha}_k$, её параметры $\boldsymbol{\theta}_k$ и множество
кластеров $I_k$, за которое он отвечает, так чтобы максимизировать
правдоподобие данных под смесью:
\begin{equation}
\label{eq:moe_likelihood}
\prod_{n=1}^{N} \sum_{k=1}^{K} p\bigl(y_n,\, m(n) \in I_k \,\big|\, x_n, \boldsymbol{\alpha}_k, \boldsymbol{\theta}_k\bigr)
\;\to\; \max_{\boldsymbol{\alpha}_{1:K},\, \boldsymbol{\theta}_{1:K},\, I_{1:K}}.
\end{equation}
Каждое слагаемое распадается на вероятность маршрутизации и условное
правдоподобие:
$p(y_n, m(n) \in I_k \mid \cdots) = p(m(n) \in I_k)\, p(y_n \mid x_n, \boldsymbol{\alpha}_k, \boldsymbol{\theta}_k, m(n) \in I_k)$,
где вероятность принадлежности есть в точности кластерная маршрутизация
$p(m(n) \in I_k) = r_{m(n), k}$.

Прямое вычисление условного множителя требует фактического обучения
сети, что вычислительно неприемлемо. Поэтому вводится
\emph{суррогатная функция}
\begin{equation}
\label{eq:surrogate_def}
u: \mathcal{A} \times [0,1]^M \to \mathbb{R}_{\geq 0},
\qquad u(\boldsymbol{\alpha}_k, \mathbf{r}_k)
\;\approx\;
- \frac{1}{|\mathcal{D}_{\mathbf{r}_k}^{\text{val}}|}
\sum_{n \in \mathcal{D}_{\mathbf{r}_k}^{\text{val}}}
\log p\bigl(y_n \mid x_n, \boldsymbol{\alpha}_k, \boldsymbol{\theta}_k^{*}\bigr),
\end{equation}
которая предсказывает \emph{среднюю по объектам кросс-энтропийную
ошибку} архитектуры $\boldsymbol{\alpha}_k$, обученной на подмножестве
кластеров, заданном вектором $\mathbf{r}_k$. Суррогат считается
постоянным внутри кластера, поэтому для каждого $n \in \mathcal{C}_m$,
приписанного эксперту $k$,
$p(y_n \mid x_n, \boldsymbol{\alpha}_k, \boldsymbol{\theta}_k, m(n) \in I_k) \approx e^{-u(\boldsymbol{\alpha}_k, \mathbf{r}_k)}$.

Подставляя эти выражения в~\eqref{eq:moe_likelihood}, логарифмируя и
группируя $N$ пообъектных слагаемых по кластерам (внутри кластера
маршрутизация и суррогат одинаковы), получаем
\emph{кластерно-зависимую целевую функцию}:
\begin{equation}
\label{eq:moe_objective}
\mathcal{J}(\boldsymbol{\alpha}_{1:K}, \mathbf{R}) \;=\;
\sum_{m=1}^{M} |\mathcal{C}_m|\,\log \sum_{k=1}^{K} r_{mk}\, e^{-u(\boldsymbol{\alpha}_k, \mathbf{r}_k)}
\;\to\; \max_{\boldsymbol{\alpha}_{1:K},\, \mathbf{R}}.
\end{equation}
Здесь $r_{mk}$ задаёт степень приписывания кластера $m$ эксперту $k$,
множитель $e^{-u(\boldsymbol{\alpha}_k, \mathbf{r}_k)}$ есть пообъектное
правдоподобие, достигаемое экспертом $k$ на своём подмножестве, а вес
$|\mathcal{C}_m|$ отражает вклад кластера, пропорциональный его размеру.
Задача~\eqref{eq:moe_objective} решается алгоритмом, описанным в
главе~\ref{ch:method}.
